\chapter{\'{E}quations canoniques}

\section{Fonction de Hamilton}\label{PAR:40_EX}

\subsection{Calcul de la fonction de Hamilton en coordonn\'{e}es cart\'{e}siennes}\label{PAR:40_EX1}

Pour un point mat\'{e}riel, nous nous basons sur le calcul de la fonction de Lagrange de la formule (\ref{EQ:4_4}) qui donne :
\benn
	L = \frac{m}{2}(\dot{x}^{2} + \dot{y}^{2} + \dot{z}^{2}) - U(x,y,z)
\eenn
Or, par d\'{e}finition des composantes de l'impulsion, voir la formule (\ref{EQ:7_5}), \`{a} savoir :
\benn
	\forall i\text{, }p_{i} = \frac{\partial L}{\partial \dot{q}_{i}}
\eenn
donc en coordonn\'{e}es cart\'{e}siennes, cela nous donne :
\bea
	p_{x} & = & \frac{\partial L}{\partial \dot{x}} = m\dot{x} \nonumber \\
	p_{y} & = & \frac{\partial L}{\partial \dot{y}} = m\dot{y} \nonumber \\
	p_{z} & = & \frac{\partial L}{\partial \dot{z}} = m\dot{z} \nonumber
\eea
La d\'{e}finition de la fonction de Hamilton donn\'{e}e par la formule (\ref{EQ:40_2}) permet d'\'{e}crire :
\bea
	H & = & m\dot{x}^{2} + m\dot{y}^{2} + m\dot{z}^{2} - \frac{m}{2}(\dot{x}^{2} + \dot{y}^{2} + \dot{z}^{2}) + U(x,y,z) = \frac{m}{2}(\dot{x}^{2} + \dot{y}^{2} + \dot{z}^{2}) + U(x,y,z) \nonumber \\
	& = & \frac{1}{2m}(p_{x}^{2} + p_{y}^{2} + p_{z}^{2}) + U(x,y,z) \nonumber
\eea

\subsection{Calcul de la fonction de Hamilton en coordonn\'{e}es cylindriques}\label{PAR:40_EX2}

En reprenant la m\^{e}me logique que dans l'exemple des coordonn\'{e}es cart\'{e}siennes mais en nous appuyant sur la d\'{e}finition de la fonction de Lagrange en coordonn\'{e}es cylindriques donn\'{e}e dans la formule (\ref{EQ:4_5}) :
\benn
	L = \frac{m}{2}(\dot{r}^{2} + r^{2}\dot{\varphi}^{2} + \dot{z}^{2}) - U(x,y,z)
\eenn
Les composantes de l'impulsion sont alors :
\bea
	p_{r} & = & \frac{\partial L}{\partial \dot{r}} = m\dot{r} \nonumber \\
	p_{\varphi} & = & \frac{\partial L}{\partial \dot{\varphi}} = mr^{2}\dot{\varphi} \nonumber \\
        p_{z} & = & \frac{\partial L}{\partial \dot{z}} = m\dot{z} \nonumber
\eea
Et la fonction de Hamilton en est d\'{e}duite :
\bea
	H & = & p_{r}\dot{r} + p_{\varphi}\dot{\varphi} + p_{z}\dot{z} - \frac{m}{2}(\dot{r}^{2} + r^{2}\dot{\varphi}^{2} + \dot{z}^{2}) + U(r,\varphi,z) \nonumber \\
	& = & m\dot{r}^{2} + mr^{2}\dot{\varphi}^{2} + m\dot{z}^{2} - \frac{m}{2}(\dot{r}^{2} + r^{2}\dot{\varphi}^{2} + \dot{z}^{2}) + U(r,\varphi,z) \nonumber \\
	& = & \frac{m}{2}(\dot{r}^{2} + r^{2}\dot{\varphi}^{2} + \dot{z}^{2}) + U(r,\varphi,z) \nonumber \\
	& = & \frac{1}{2m}\left(p_{r}^{2} + \frac{p_{\varphi}}{r^{2}}i + p_{z}\right)  + U(r,\varphi,z) \nonumber
\eea

\subsection{Calcul de la fonction de Hamilton en coordonn\'{e}es sph\'{e}riques}\label{PAR:40_EX3}

Nous allons nous baser ici sur la fonction de Lagrange en coordonn\'{e}es sph\'{e}riques de la formule (\ref{EQ:4_6}) :
\benn
	L = \frac{m}{2}(\dot{r}^{2} + r^{2}\dot{\theta}^{2} + r^{2}\sin^{2}\theta\dot{\varphi}^{2}) - U(x,y,z)
\eenn
Les composantes de l'impulsion dans ce syst\`{e}me de coordonn\'{e}es sont :
\bea
        p_{r} & = & \frac{\partial L}{\partial \dot{r}} = m\dot{r} \nonumber \\
        p_{\theta} & = & \frac{\partial L}{\partial \dot{\theta}} = mr^{2}\dot{\theta} \nonumber \\
	p_{\varphi} & = & \frac{\partial L}{\partial \dot{\varphi}} = mr^{2}\sin^{2}\theta\dot{\varphi} \nonumber
\eea
Et la fonction de Hamilton se calcule ainsi :
\bea
	H & = & _{r}\dot{r} + p_{\theta}\dot{\theta} + p_{\varphi}\dot{\varphi} - \frac{m}{2}(\dot{r}^{2} + r^{2}\dot{\theta}^{2} + r^{2}\sin^{2}\theta\dot{\varphi}^{2}) + U(r,\theta,\varphi) \nonumber \\
	& = & m\dot{r}^{2} + mr^{2}\dot{\theta}^{2} + mr^{2}\sin^{2}\theta\dot{\varphi}^{2} - \frac{m}{2}(\dot{r}^{2} + r^{2}\dot{\theta}^{2} + r^{2}\sin^{2}\theta\dot{\varphi}^{2}) + U(r,\theta,\varphi) \nonumber \\
	& = & \frac{m}{2}(\dot{r}^{2} + r^{2}\dot{\theta}^{2} + r^{2}\sin^{2}\theta\dot{\varphi}^{2}) + U(r,\theta,\varphi) \nonumber \\
	& = & \frac{1}{2m}\left(p_{r}^{2} + \frac{p_{\theta}^{2}}{r^{2}} + \frac{p_{\varphi}^{2}}{r^{2}\sin^{2}\theta}\right) + U(r,\theta,\varphi) \nonumber
\eea

\subsection{Cas d'un syst\`{e}me anim\'{e} d'une rotation uniforme}\label{PAR:40_EX4}

\subsection{Cas d'un syst\`{e}me \`{a} $n+1$ particules}\label{PAR:40_EX5}

Nous allons ici calculer la fonction de Hamilton dans le cas o\`{u} le syst\`{e}me est compos\'{e} d'une particule de masse $M$ et de $n$ particules de masses $m$ sachant que le mouvement est rapport\'{e} au r\'{e}f\'{e}rentiel du centre d'inertie. Nous nous r\'{e}f\'{e}rons alors \`{a} la fonction de Lagrange calcul\'{e}e dans le paragraphe (\ref{PAR:13}) qui donne :
\bea
	L & = & \frac{m}{2}\sum_{a}\dot{\vec{r}}_{a}^{\,2} - \frac{m^{2}}{2(M + nm)}\left(\sum_{a}\dot{\vec{r}}_{a}\right)^{2} - U \nonumber \\
	& = & \frac{m}{2}\sum_{a}\vec{v}_{a}^{\,2} - \frac{m^{2}}{2(M + nm)}\left(\sum_{a}\vec{v}_{a}\right)^{2} - U \nonumber
\eea
L'impulsion de la particule $a$ a pour d\'{e}finition :
\benn
	\vec{p}_{a} = \frac{\partial L}{\partial \vec{v}_{a}}
\eenn
Et en notant que :
\benn
	\frac{\partial(\sum_{a}\vec{v}_{a})^{2}}{\partial\vec{v}_{a}} = 2\sum_{a}\vec{v}_{a}\times\frac{\partial\vec{v}_{a}}{\partial\vec{v}_{a}} = 2\sum_{a}\vec{v}_{a}
\eenn
alors l'impulsion se d\'{e}veloppe ainsi :
\benn
	\vec{p}_{a} = m\vec{v}_{a} - \frac{m^{2}}{M + nm}\sum_{a}\vec{v}_{a} \Rightarrow \sum_{a}\vec{p}_{a} = m\sum_{a}\vec{v}_{a} - \frac{nm^{2}}{M + nm}\sum_{a}\vec{v}_{a}
\eenn
car $\sum_{a}(\sum_{a}\vec{v}_{a}) = n\sum_{a}\vec{v}_{a}$. Nous en d\'{e}duisons donc :
\benn
	\sum_{a}\vec{p}_{a} = \frac{m}{M + nm}(M + nm - nm)\sum_{a}\vec{v}_{a} = \frac{mM}{M + nm}\sum_{a}\vec{v}_{a}
\eenn
et :
\benn
	\vec{p}_{a} = m\vec{v}_{a} - \frac{m^{2}}{M + nm}\sum_{a}\vec{v}_{a} = m\vec{v}_{a} - \frac{m}{M}\sum_{a}\vec{p}_{a} \Leftrightarrow \vec{v}_{a} = \frac{\vec{p}_{a}}{m} + \frac{\sum_{a}\vec{p}_{a}}{M}
\eenn
\`{A} partir de tous ces \'{e}l\'{e}ments, la fonction de Lagrange devient alors en l'\'{e}crivant en fonction des impulsions :
\bea
	L & = & \frac{m}{2}\sum_{a}\left(\frac{\vec{p}_{a}}{m} + \frac{\sum_{a}\vec{p}_{a}}{M}\right)^{2} - \frac{m^{2}}{2(M + nm)}\left(\frac{M + nm}{mM}\sum_{a}\vec{p}_{a}\right)^{2} - U \nonumber \\
	& = & \frac{1}{2m}\sum_{a}\vec{p}_{a}^{\,2} + \frac{m}{2M^{2}}\sum_{a}(\sum_{a}\vec{p}_{a})^{2} + \frac{1}{M}\sum_{a}(\vec{p}_{a}\sum_{a}\vec{p}_{a}) - \frac{M + nm}{2M^{2}}(\sum_{a}\vec{p}_{a})^{2} - U \nonumber \\
	& = & \frac{1}{2m}\sum_{a}\vec{p}_{a}^{\,2} - \frac{1}{2M}(\sum_{a}\vec{p}_{a})^{2} + \frac{1}{M}\sum_{a}(\vec{p}_{a}\sum_{a}\vec{p}_{a}) - U \nonumber
\eea
ce qui nous permet de calculer la fonction de Hamilton du syst\`{e}me :
\bea
	H & = & \sum_{a}\vec{p}_{a}\vec{v}_{a} - L \nonumber \\
	& = & \sum_{a}\vec{p}_{a}\left(\frac{\vec{p}_{a}}{m} + \frac{\sum_{a}\vec{p}_{a}}{M}\right) - \frac{1}{2m}\sum_{a}\vec{p}_{a}^{\,2} + \frac{1}{2M}(\sum_{a}\vec{p}_{a})^{2} - \frac{1}{M}\sum_{a}(\vec{p}_{a}\sum_{a}\vec{p}_{a}) + U \nonumber \\
	& = & \frac{1}{m}\sum_{a}\vec{p}_{a}^{\,2} + \frac{1}{M}\sum_{a}\left(\vec{p}_{a}\sum_{a}\vec{p}_{a}\right) - \frac{1}{2m}\sum_{a}\vec{p}_{a}^{\,2} + \frac{1}{2M}\left(\sum_{a}\vec{p}_{a}\right)^{2} - \frac{1}{M}\sum_{a}\left(\vec{p}_{a}\sum_{a}\vec{p}_{a}\right) + U \nonumber \\
	& = & \frac{1}{2m}\sum_{a}\vec{p}_{a}^{\,2} + \frac{1}{2M}\left(\sum_{a}\vec{p}_{a}\right)^{2} + U \nonumber
\eea

\section{Fonction de Routh}\label{PAR:41_EX}

Il s'agit ici de trouver la fonction de Routh dans le cadre du mouvement d'une toupie sym\'{e}trique dans un champ ext\'{e}rieur $U(\varphi,\theta)$ en \'{e}liminant la coordonn\'{e}e cyclique $\psi$ sachant que $\varphi$, $\theta$ et $\psi$ sont les angles d'Euler qui ont \'{e}t\'{e} \'{e}tudi\'{e}s dans le paragraphe (\ref{PAR:35}).

La fonction de Lagrange est donn\'{e}e par la r\'{e}solution de l'exercice (\ref{PAR:35_EX1}) :
\benn
	L = \frac{I'_{1}}{2}(\dot{\theta}^{2} + \sin^{2}\theta\dot{\varphi}^{2}) + \frac{I_{3}}{2}(\dot{\psi} + \cos\theta\dot{\varphi})^{2} - U(\theta,\varphi)
\eenn
La d\'{e}finition de la fonction de Routh permet d\'{e}crire pour notre cas :
\benn
	R = p_{\theta}\dot{\theta} + p_{\varphi}\dot{\varphi} + p_{\psi}\dot{\psi} - L
\eenn
Les impulsions se calculent ainsi :
\bea
	p_{\theta} & = & \frac{\partial L}{\partial\dot{q}_{\theta}} = I'_{1}\dot{\theta} \nonumber \\
	p_{\varphi} & = & \frac{\partial L}{\partial\dot{q}_{\varphi}} = I'_{1}\sin^{2}\theta\dot{\varphi} + I_{3}\sin^{2}\theta\dot{\varphi} + I_{3}\cos\theta\dot{\psi} \nonumber \\
	p_{\psi} & = & \frac{\partial L}{\partial\dot{q}_{\psi}} = I_{3}\dot{\psi} + I_{3}\cos\theta\dot{\varphi} \nonumber
\eea
La coordonn\'{e}e $\psi$ est cylique donc $\frac{\partial L}{\partial q_{\psi}} = 0$ et $p_{\psi} = \frac{\partial L}{\partial\dot{q}_{\psi}}$ a une valeur constante, ce qui permet d'en d\'{e}duire :
\benn
	\dot{\psi} = \frac{p_{\psi}}{I_{3}} - \cos\theta\dot{\varphi}
\eenn
Le calcul de la fonction de Routh du syst\`{e}me se r\'{e}alise ainsi :
\bea
	R & = & I'_{1}\dot{\theta}^{2} + I'_{1}\sin^{2}\theta\dot{\varphi}^{2} + I_{3}\cos^{2}\theta\dot{\varphi}^{2} + I_{3}\cos\theta\dot{\varphi}\dot{\psi} + p_{\psi}\dot{\psi} \nonumber \\
	& & - \frac{I'_{1}}{2}\dot{\theta}^{2} - \frac{I'_{1}}{2}\sin^{2}\theta\dot{\varphi}^{2} - \frac{I_{3}}{2}\dot{\psi}^{2} - I_{3}\cos\theta\dot{\varphi}\dot{\psi} - \frac{I_{3}}{2}\cos^{2}\theta\dot{\varphi}^{2} + U \nonumber \\
	& = & \frac{I'_{1}}{2}\dot{\theta}^{2} + \frac{I'_{1}}{2}\sin^{2}\theta\dot{\varphi}^{2} + \frac{I_{3}}{2}\cos^{2}\theta\dot{\varphi}^{2} + p_{\psi}\dot{\psi} - \frac{I_{3}}{2}\dot{\psi}^{2} + U \nonumber \\
	& = & \frac{I'_{1}}{2}(\dot{\theta}^{2} + \sin^{2}\theta\dot{\varphi}^{2}) + \frac{I_{3}}{2}\cos^{2}\theta\dot{\varphi}^{2} + p_{\psi}\left(\frac{p_{\psi}}{I_{3}} - \cos\theta\dot{\varphi}\right) - \frac{I_{3}}{2}\left(\frac{p_{\psi}^{2}}{I_{3}^{2}} + \cos^{2}\theta\dot{\varphi}^{2} - 2\frac{p_{\psi}}{I_{3}}\cos\theta\dot{\varphi}\right) + U \nonumber \\
	\Leftrightarrow R & = & \frac{I'_{1}}{2}(\dot{\theta}^{2} + \sin^{2}\theta\dot{\varphi}^{2}) + \frac{p_{\psi}^{2}}{2I_{3}} + U \nonumber
\eea
Nous arrivons ici \`{a} une fonction de Routh bien diff\'{e}rente de celle de la solution dans le livre de Landau ...

\section{Crochets de Poisson}\label{PAR:42_EX}

\subsection{Impulsion et moment cin\'{e}tique}\label{PAR:42_EX1}

Nous allons ici d\'{e}terminer les crochets de Poisson pour l'impulsion $\vec{p}$ et le moment cin\'{e}tique $\vec{M}$ en composantes cart\'{e}siennes. Par d\'{e}finition, le moment cin\'{e}tique vaut :
\benn
	\vec{M} = \vec{r}\wedge\vec{p} = \begin{pmatrix} yp_{z} - zp_{y} \\ zp_{x} - xp_{z} \\ xp_{y} - yp_{x} \end{pmatrix}
\eenn
donc :
\bea
	\{M_{x},p_{x}\} & = & \frac{\partial M_{x}}{\partial p_{x}}\frac{\partial p_{x}}{\partial x} - \frac{\partial M_{x}}{\partial x}\frac{\partial p_{x}}{\partial p_{x}} + \frac{\partial M_{x}}{\partial p_{y}}\frac{\partial p_{x}}{\partial y} - \frac{\partial M_{x}}{\partial y}\frac{\partial p_{x}}{\partial p_{y}} + \frac{\partial M_{x}}{\partial p_{z}}\frac{\partial p_{x}}{\partial z} - \frac{\partial M_{x}}{\partial z}\frac{\partial p_{x}}{\partial p_{z}} \nonumber \\
	& = & -\frac{\partial M_{x}}{\partial x} = 0 \nonumber \\
	\{M_{x},p_{y}\} & = & \frac{\partial M_{x}}{\partial p_{x}}\frac{\partial p_{y}}{\partial x} - \frac{\partial M_{x}}{\partial x}\frac{\partial p_{y}}{\partial p_{x}} + \frac{\partial M_{x}}{\partial p_{y}}\frac{\partial p_{y}}{\partial y} - \frac{\partial M_{x}}{\partial y}\frac{\partial p_{y}}{\partial p_{y}} + \frac{\partial M_{x}}{\partial p_{z}}\frac{\partial p_{y}}{\partial z} - \frac{\partial M_{x}}{\partial z}\frac{\partial p_{y}}{\partial p_{z}} \nonumber \\
	& = & -\frac{\partial M_{x}}{\partial y} = -p_{z} \nonumber \\
	\{M_{x},p_{z}\} & = & \frac{\partial M_{x}}{\partial p_{x}}\frac{\partial p_{z}}{\partial x} - \frac{\partial M_{x}}{\partial x}\frac{\partial p_{z}}{\partial p_{x}} + \frac{\partial M_{x}}{\partial p_{y}}\frac{\partial p_{z}}{\partial y} - \frac{\partial M_{x}}{\partial y}\frac{\partial p_{z}}{\partial p_{y}} + \frac{\partial M_{x}}{\partial p_{z}}\frac{\partial p_{z}}{\partial z} - \frac{\partial M_{x}}{\partial z}\frac{\partial p_{z}}{\partial p_{z}} \nonumber \\ 
        & = & -\frac{\partial M_{x}}{\partial z} = 0 \nonumber
\eea
Par rotation des indices, nous arrivons \`{a} :
\bea
	\{M_{y},p_{x}\} & = & -\frac{\partial M_{y}}{\partial x} = p_{z}\text{, }\{M_{y},p_{y}\} = -\frac{\partial M_{y}}{\partial y} = 0\text{ et }\{M_{y},p_{z}\} = -\frac{\partial M_{y}}{\partial z} = -p_{x} \nonumber \\
	\{M_{z},p_{x}\} & = & -\frac{\partial M_{z}}{\partial x} = p_{y}\text{, }\{M_{z},p_{y}\} = -\frac{\partial M_{z}}{\partial y} = -p_{x}\text{ et }\{M_{z},p_{z}\} = -\frac{\partial M_{z}}{\partial z} = 0 \nonumber
\eea

\subsection{Crochets de Poisson des composantes du moment cin\'{e}tique}\label{PAR:42_EX2}

Par application de la d\'{e}fintion des crochets de Poisson, voir (\ref{EQ:42_1}) :
\bea
	\{M_{x},M_{x}\} & = & 0 \nonumber \\
	\{M_{x},M_{y}\} & = & \frac{\partial M_{x}}{\partial p_{x}}\frac{\partial M_{y}}{\partial x} - \frac{\partial M_{x}}{\partial x}\frac{\partial M_{y}}{\partial p_{x}} + \frac{\partial M_{x}}{\partial p_{y}}\frac{\partial M_{y}}{\partial y} - \frac{\partial M_{x}}{\partial y}\frac{\partial M_{y}}{\partial p_{y}} + \frac{\partial M_{x}}{\partial p_{z}}\frac{\partial M_{y}}{\partial z} - \frac{\partial M_{x}}{\partial z}\frac{\partial M_{y}}{\partial p_{z}} \nonumber \\
	& = & 0 - 0 + 0 - 0 + yp_{x} - (-p_{y})(-x) \nonumber \\
        & = & -M_{z} \nonumber \\
	\{M_{y},M_{z}\} & = & \frac{\partial M_{y}}{\partial p_{x}}\frac{\partial M_{z}}{\partial x} - \frac{\partial M_{y}}{\partial x}\frac{\partial M_{z}}{\partial p_{x}} + \frac{\partial M_{y}}{\partial p_{y}}\frac{\partial M_{z}}{\partial y} - \frac{\partial M_{y}}{\partial y}\frac{\partial M_{z}}{\partial p_{y}} + \frac{\partial M_{y}}{\partial p_{z}}\frac{\partial M_{z}}{\partial z} - \frac{\partial M_{y}}{\partial z}\frac{\partial M_{y}}{\partial p_{z}} \nonumber \\
	& = & zp_{y} - (-p_{z})(-y) + 0 - 0 + 0 - 0 \nonumber \\
        & = & -M_{x} \nonumber \\
	\{M_{z},M_{x}\} & = & \frac{\partial M_{z}}{\partial p_{x}}\frac{\partial M_{x}}{\partial x} - \frac{\partial M_{z}}{\partial x}\frac{\partial M_{x}}{\partial p_{x}} + \frac{\partial M_{z}}{\partial p_{y}}\frac{\partial M_{x}}{\partial y} - \frac{\partial M_{z}}{\partial y}\frac{\partial M_{x}}{\partial p_{y}} + \frac{\partial M_{z}}{\partial p_{z}}\frac{\partial M_{x}}{\partial z} - \frac{\partial M_{z}}{\partial z}\frac{\partial M_{x}}{\partial p_{z}} \nonumber \\
	& = & 0 - 0 + xp_{z} - (-p_{x})(-z) \nonumber \\
        & = & -M_{y} \nonumber
\eea

\subsection{Crochets de Poisson entre une fonction scalaire et le moment cin\'{e}tique}\label{PAR:42_EX3}

\subsection{Crochets de Poisson entre une fonction vectorielle et le moment cin\'{e}tique}\label{PAR:42_EX4}

Soit $\vec{f}$ une fonction vectorielle des coordonn\'{e}es et de l'impulsion d'une particule et $\vec{i}$, $\vec{j}$ et $\vec{k}$ les vecteurs unitaires respectivement pour les axes $x$, $y$ et $z$. La fonction $\vec{f}$ peut se d\'{e}composer de cette mani\`{e}re dans le r\'{e}f\'{e}rentiel des coordonn\'{e}es cart\'{e}siennes car $\vec{r}$, $\vec{p}$ et $\vec{r}\wedge\vec{p}$ sont orthogonaux entre eux :
\benn
	\vec{f} = f_{1}\vec{r} + f_{2}\vec{p} + f_{3}\vec{r}\wedge\vec{p}
\eenn
avec $f_{1}$, $f_{2}$ et $f_{3}$ trois fonctions scalaires des coordonn\'{e}es et de l'impulsion.

Nous rappelons que :
\begin{itemize}
	\item le moment cin\'{e}tique $\vec{M}$ vaut exactement $\vec{r}\wedge\vec{p}$, soit :
\benn
	\vec{M} = \begin{pmatrix} yp_{z} - zp_{y} \\ zp_{x} - xp_{z} \\ xp_{y} - yp_{x}  \end{pmatrix}
\eenn
	\item en conclusion de l'exercice pr\'{e}c\'{e}dent (\ref{PAR:42_EX2}), soit $f$ une fonction scalaire des coordonn\'{e}es et de l'impulsion, alors $\{f,M_{z}\} = 0$
\end{itemize}
Nous pouvons donc d\'{e}velopper la quantit\'{e} en appliquant la lin\'{e}arit\'{e} des crochets de Poisson, voir (\ref{EQ:42_8}) puis la formule (\ref{EQ:42_9}) :
\bea
	\{\vec{f},M_{z}\} & = & \{f_{1}\vec{r},M_{z}\} + \{f_{2}\vec{p},M_{z}\} + \{f_{3}\vec{r}\wedge\vec{p},M_{z}\} \nonumber \\
	& = & f_{1}\{\vec{r},M_{z}\} + \vec{r}\{f_{1},M_{z}\} + f_{2}\{\vec{p},M_{z}\} + \vec{p}\{f_{2},M_{z}\} + f_{3}\{\vec{r}\wedge\vec{p},M_{z}\} + \vec{r}\wedge\vec{p}\{f_{3},M_{z}\} \nonumber
\eea
Et en suivant le second rappel plus haut, soit $\{f_{1},M_{z}\} = \{f_{2},M_{z}\} = \{f_{3},M_{z}\} = 0$, nous avons :
\benn
	\{\vec{f},M_{z}\} = f_{1}\{\vec{r},M_{z}\} + f_{2}\{\vec{p},M_{z}\} + f_{3}\{\vec{r}\wedge\vec{p},M_{z}\}
\eenn
En appliquant la formule (\ref{EQ:42_11}), nous pouvons \'{e}crire :
\benn
	\{\vec{r},M_{z}\} = -\frac{\partial M_{z}}{\partial\vec{p}} = -\frac{\partial M_{z}}{\partial p_{x}}\vec{i} - \frac{\partial M_{z}}{\partial p_{y}}\vec{j} - \frac{\partial M_{z}}{\partial p_{z}}\vec{k} = y\vec{i} - x\vec{j}
\eenn
En appliquant, cette fois-ci, la formule (\ref{EQ:42_12}) :
\benn
	\{\vec{p},M_{z}\} = \frac{\partial M_{z}}{\partial\vec{r}} = \frac{\partial M_{z}}{\partial x}\vec{i} + \frac{\partial M_{z}}{\partial y}\vec{j} + \frac{\partial M_{z}}{\partial z}\vec{k} = p_{y}\vec{i} - p_{x}\vec{j}
\eenn
Nous devons maintenant calculer la quantit\'{e} :
\benn
	\{\vec{M},M_{z}\} = \frac{\partial\vec{M}}{\partial p_{x}}\frac{\partial M_{z}}{\partial x} - \frac{\partial\vec{M}}{\partial x}\frac{\partial M_{z}}{\partial p_{x}} + \frac{\partial\vec{M}}{\partial p_{y}}\frac{\partial M_{z}}{\partial y} - \frac{\partial\vec{M}}{\partial y}\frac{\partial M_{z}}{\partial p_{y}} + \frac{\partial\vec{M}}{\partial p_{z}}\frac{\partial M_{z}}{\partial z} - \frac{\partial\vec{M}}{\partial z}\frac{\partial M_{z}}{\partial p_{z}}
\eenn
sachant qu'avec la d\'{e}composition de $\vec{M}$, nous pouvons en d\'{e}duire :
\bea
	\frac{\partial\vec{M}}{\partial x} & = & -p_{z}\vec{j} + p_{y}\vec{k} \text{, } \frac{\partial\vec{M}}{\partial y} = p_{z}\vec{i} - p_{x}\vec{k} \text{ et } \frac{\partial\vec{M}}{\partial z} = -p_{y}\vec{i} + p_{x}\vec{j} \nonumber \\
	\frac{\partial\vec{M}}{\partial p_{x}} & = & z\vec{j} - y\vec{k} \text{, } \frac{\partial\vec{M}}{\partial p_{y}} = -z\vec{i} + x\vec{k} \text{ et } \frac{\partial\vec{M}}{\partial p_{z}} = y\vec{i} - x\vec{j} \nonumber
\eea

\section{Principe de Maupertuis}\label{PAR:44_EX}

L'objectif est ici de d\'{e}terminer l'\'{e}quation diff\'{e}rentielle de la trajectoire \`{a} partir du principe variationnel trouv\'{e} en (\ref{EQ:44_10}) : $\delta\left(\int{\sqrt{2m(E - U)}\mathrm{d}l}\right) = 0$. D\'{e}veloppons la variation :
\benn
	\delta\left(\int{\sqrt{2m(E - U)}\mathrm{d}l}\right) = \int{\frac{\partial}{\partial\vec{r}}(\sqrt{2m(E - U)}\mathrm{d}l)\delta\vec{r}} = \int{\frac{\partial}{\partial\vec{r}}(\sqrt{2m(E - U)})\mathrm{d}l\delta\vec{r}} + \int{\sqrt{2m(E - U)}\frac{\partial\mathrm{d}l}{\partial\vec{r}}\delta\vec{r}}
\eenn
Parce que $E$ est constante et $U$ est une fonction uniquement de la position $\vec{r}$ alors :
\benn
	\int{\frac{\partial}{\partial\vec{r}}(\sqrt{2m(E - U))}\mathrm{d}l\delta\vec{r}} = -\int{\frac{\sqrt{m}\delta\vec{r}}{\sqrt{2(E - U)}}\frac{\partial U}{\partial\vec{r}}\mathrm{d}l}
\eenn
Par d\'{e}finition de l'\'{e}l\'{e}ment de longueur le long de la trajectoire : $\mathrm{d}l^{2} = \mathrm{d}\vec{r}^{\,2}$ donc :
\bea
	\delta(\mathrm{d}l^{2}) & = & \delta(\mathrm{d}\vec{r}^{\,2}) \Leftrightarrow 2\mathrm{d}l\delta(\mathrm{d}l) = 2\mathrm{d}\vec{r}\delta(\mathrm{d}\vec{r}) = 2\mathrm{d}\vec{r}\mathrm{d}(\delta\vec{r}) \nonumber \\
	\Leftrightarrow \frac{\mathrm{d}\vec{r}}{\mathrm{d}l}\mathrm{d}(\delta\vec{r}) & = & \delta(\mathrm{d}l) = \mathrm{d}(\delta l) = \frac{\partial\delta l}{\partial\vec{r}}\mathrm{d}\vec{r} \nonumber
\eea
Donc la variation recherch\'{e}e se d\'{e}veloppe ainsi :
\benn
	\delta\left(\int{\sqrt{2m(E - U)}\mathrm{d}l}\right) = -\int{\frac{\sqrt{m}\delta\vec{r}}{\sqrt{2(E - U)}}\frac{\partial U}{\partial\vec{r}}\mathrm{d}l} + \int{\sqrt{2m(E - U)}\frac{\mathrm{d}\vec{r}}{\mathrm{d}l}\mathrm{d}(\delta\vec{r})}
\eenn
Nous pouvons utiliser l'int\'{e}gration par partie pour le second terme tel que :
\benn
	\int{\sqrt{2m(E - U)}\frac{\mathrm{d}\vec{r}}{\mathrm{d}l}\mathrm{d}(\delta\vec{r})} = \left[\sqrt{2m(E - U)}\mathrm{d}\vec{r}\delta\vec{r}\right] - \int{\frac{\mathrm{d}}{\mathrm{d}l}\left(\sqrt{2m(E - U)}\frac{\mathrm{d}\vec{r}}{\mathrm{d}l}\right)\delta\vec{r}\mathrm{d}l}
\eenn
Nous choisissons les conditions aux limites telles que le premier terme du second membre soit nul. Donc l'\'{e}galit\'{e} (\ref{EQ:44_10}) permet d'\'{e}crire :
\bea
	\delta\left(\int{\sqrt{2m(E - U)}\mathrm{d}l}\right) & = & 0 \nonumber \\
	\Leftrightarrow -\int{\frac{\sqrt{m}\delta\vec{r}}{\sqrt{2(E - U)}}\frac{\partial U}{\partial\vec{r}}\mathrm{d}l} - \int{\frac{\mathrm{d}}{\mathrm{d}l}\left(\sqrt{2m(E - U)}\frac{\mathrm{d}\vec{r}}{\mathrm{d}l}\right)\delta\vec{r}\mathrm{d}l} &= & 0 \nonumber \\
	\Leftrightarrow \sqrt{2m}\frac{\mathrm{d}}{\mathrm{d}l}\left(\sqrt{E - U}\frac{\mathrm{d}\vec{r}}{\mathrm{d}l}\right) & = & -\frac{\sqrt{m}}{2}\frac{\partial U}{\partial\vec{r}}\frac{1}{\sqrt{E - U}} \nonumber \\
	\Leftrightarrow 2\sqrt{E - U}\frac{\mathrm{d}}{\mathrm{d}l}\left(\sqrt{E - U}\frac{\mathrm{d}\vec{r}}{\mathrm{d}l}\right) & = & -\frac{\partial U}{\partial\vec{r}} \nonumber
\eea
Nous pouvons remarquer que cette \'{e}quation ne d\'{e}pend pas de la masse de la particule en mouvement. Par d\'{e}finition, la quantit\'{e} $-\frac{\partial U}{\partial\vec{r}}$ est une force que nous appelons $\vec{F}$. Ainsi, nous \'{e}crivons :
\bea
	\frac{\mathrm{d}}{\mathrm{d}l}\left(\sqrt{E - U}\frac{\mathrm{d}\vec{r}}{\mathrm{d}l}\right) & = & \frac{\mathrm{d}}{\mathrm{d}l}\left(\sqrt{E - U}\right)\frac{\mathrm{d}\vec{r}}{\mathrm{d}l} + \sqrt{E - U}\frac{\mathrm{d}^{2}\vec{r}}{\mathrm{d}l^{2}} \nonumber \\
	& = & \left(\frac{\mathrm{d}}{\mathrm{d}\vec{r}}\left(\sqrt{E - U}\right)\cdot\frac{\mathrm{d}\vec{r}}{\mathrm{d}l}\right)\cdot\frac{\mathrm{d}\vec{r}}{\mathrm{d}l} + \sqrt{E - U}\frac{\mathrm{d}^{2}\vec{r}}{\mathrm{d}l^{2}} \nonumber \\
	& = & -\frac{\frac{\partial U}{\partial\vec{r}}\cdot\frac{\mathrm{d}\vec{r}}{\mathrm{d}l}}{2\sqrt{E - U}}\cdot\frac{\mathrm{d}\vec{r}}{\mathrm{d}l} + \sqrt{E - U}\frac{\mathrm{d}^{2}\vec{r}}{\mathrm{d}l^{2}} \nonumber
\eea
En posant $\vec{u}_{r} = \frac{\mathrm{d}\vec{r}}{\mathrm{d}l}$ le vecteur unitaire tangentiel \`{a} la trajectoire, alors l'\'{e}quation pr\'{e}c\'{e}dente devient :
\bea
	(\vec{F}\cdot\vec{u}_{r})\cdot\vec{u}_{r} + 2(E - U)\frac{\mathrm{d}^{2}\vec{r}}{\mathrm{d}l^{2}} & = & \vec{F} \nonumber \\
	\Leftrightarrow \frac{\mathrm{d}^{2}\vec{r}}{\mathrm{d}l^{2}} & = & \frac{\vec{F} - (\vec{F}\cdot\vec{u}_{r})\cdot\vec{u}_{r}}{2(E - U)} \nonumber
\eea
o\`{u} $\vec{F} - (\vec{F}\cdot\vec{u}_{r})\cdot\vec{u}_{r}$ est la force \`{a} laquelle il est retir\'{e}e sa composante tangentielle, il s'agit donc de sa composante normale, orthogonale, \`{a} la trajectoire, soit $\vec{F}_{n}$. La g\'{e}om\'{e}trie diff\'{e}rentielle donne $\frac{\mathrm{d}^{2}\vec{r}}{\mathrm{d}l^{2}} = \frac{\vec{n}}{R}$ avec $\vec{n}$ le vecteur unitaire normal et $R$ le rayon de courbure de la trajectoire. Alors l'\'{e}quation devient :
\benn
	\frac{\vec{n}}{R} = \frac{\vec{F}_{n}}{2(E - U)}
\eenn
Mais par d\'{e}finition de l'\'{e}nergie totale du syst\`{e}me $E = T + U$, soit $E - U = T = \frac{1}{2}m\vec{v}^{\,2}$, ce qui donne finalement :
\benn
	\frac{m\vec{v}^{\,2}}{R}\vec{n} = \vec{F}_{n}
\eenn
qui correspond bien \`{a} l'acc\'{e}l\'{e}ration normale pour un mouvement le long d'une trajectoire courbe.

\section{S\'{e}paration des variables}\label{PAR:48_EX}

\subsection{Superposition d'un champ coulombien et d'un champ homog\`{e}ne}\label{PAR:48_EX1}

Il s'agit ici de trouver l'int\'{e}grale compl\`{e}te de l'\'{e}quation de Hamilton-Jacobi pour une particule de masse $m$ en mouvement dans un champ compos\'{e} d'un champ coulombien et d'un champ homog\`{e}ne tel que :
\benn
	U = \frac{\alpha}{r} - Fz
\eenn
En utilisant les coordonn\'{e}es paraboliques (\ref{EQ:48_10}) et (\ref{EQ:48_11}), nous pouvons r\'{e}\'{e}crire l'\'{e}nergie potentielle telle que :
\benn
	U(\xi,\eta) = \frac{2\alpha}{\xi + \eta} - \frac{F}{2}(\xi - \eta) = \frac{\alpha - \frac{F}{2}\xi^{2} + \alpha + \frac{F}{2}\eta^{2}}{\xi +\eta}
\eenn
En identifiant cette relation avec (\ref{EQ:48_15}), nous obtenons :
\benn
	a(\xi) = \alpha - \frac{F}{2}\xi^{2}\text{ et }b(\eta) = \alpha + \frac{F}{2}\eta^{2}
\eenn
L'utilisation de la relation (\ref{EQ:48_16}) nous donne de fait l'int\'{e}grale compl\`{e}te solution de l'\'{e}quation de Hamilton-Jacobi. Les \'{e}quations issues de la m\'{e}thode de s\'{e}paration des variables sont :
\bea
	2\xi p_{\xi}^{2} + ma(\xi) - m\xi E +\frac{p_{\varphi}^{2}}{2\xi} & = & \beta \nonumber \\
	2\eta p_{\eta}^{2} + mb(\eta) - m\eta E + \frac{p_{\varphi}^{2}}{2\eta} & = & -\beta \nonumber
\eea

La conservation de l'\'{e}nergie nous apporte :
\benn
	H = E \Leftrightarrow \frac{2}{m}\frac{\xi p_{\xi}^{2} + \eta p_{\eta}^{2}}{\xi + \eta} + \frac{p_{\varphi}^{2}}{2m\xi\eta} + \frac{\alpha}{r} - Fz = E
\eenn
alors que la diff\'{e}rence des deux \'{e}quations issues de la s\'{e}paration des variables va nous permettre d'acc\'{e}der \`{a} la quantit\'{e} $\beta$ :
\bea
	2\beta & = & 2\xi p_{\xi}^{2} - 2\eta p_{\eta}^{2} + m(a(\xi) - b(\eta)) + mE(\eta - \xi) + \frac{p_{\varphi}^{2}}{2}\left(\frac{1}{\xi} - \frac{1}{\eta}\right) \nonumber \\
	& = & 2\xi p_{\xi}^{2} - 2\eta p_{\eta}^{2} - \frac{mF}{2}(\xi^{2} + \eta^{2}) + \frac{p_{\varphi}^{2}}{2}\frac{(\eta - \xi)}{\xi\eta} + m(\eta - \xi)\left(\frac{2}{m}\frac{\xi p_{\xi}^{2} + \eta p_{\eta}^{2}}{\xi + \eta} + \frac{p_{\varphi}^{2}}{2m\xi\eta} + \frac{\alpha}{r} - Fz\right) \nonumber \\
	& = & 2\xi p_{\xi}^{2} - 2\eta p_{\eta}^{2} - \frac{mF}{2}(\xi^{2} + \eta^{2}) + \frac{p_{\varphi}^{2}}{2}\frac{(\eta - \xi)}{\xi\eta} \nonumber \\
	& & -2\frac{(\xi - \eta)}{(\xi + \eta)}(\xi p_{\xi}^{2} + \eta p_{\eta}^{2}) - \frac{(\xi - \eta)}{2\xi\eta}p_{\varphi}^{2} - m(\xi - \eta)\left(\frac{\alpha}{r} - Fz\right) \nonumber \\
	& = & 2\xi p_{\xi}^{2}\left(1 - \frac{\xi - \eta}{\xi + \eta}\right) - 2\eta p_{\eta}^{2}\left(1 + \frac{\xi - \eta}{\xi + \eta}\right) - \frac{mF}{2}(\xi^{2} + \eta^{2}) + mFz(\xi - \eta) \nonumber \\
	& & -m(\xi - \eta)\frac{\alpha}{r} - \frac{\xi - \eta}{\xi\eta}p_{\varphi}^{2} \nonumber
\eea
Or nous calculons les facteurs suivants de l'\'{e}quation pr\'{e}c\'{e}dente :
\bea
	1 - \frac{\xi - \eta}{\xi + \eta} & = & 2\frac{\eta}{\xi + \eta} \nonumber \\
	1 + \frac{\xi - \eta}{\xi + \eta} & = & 2\frac{\xi}{\xi + \eta} \nonumber \\
	-\frac{mF}{2}(\xi^{2} + \eta^{2}) + mFz(\xi - \eta) & = & -\frac{mF}{2}\left(2(\rho^{2} + 2z^{2}) - 4z^{2})\right) = -m\rho^{2}F \nonumber \\
	-m(\xi - \eta)\frac{\alpha}{r} & = & -2m\alpha\frac{z}{\rho} \nonumber \\
	-\frac{\xi - \eta}{\xi\eta} & = & -2\frac{z}{\rho^{2}} \nonumber
\eea
Le passage aux coordonn\'{e}es $\rho$ et $z$ demande les calculs suivants des impulsions sachant que $\rho = \sqrt{\xi\eta}$ et $z = \frac{1}{2}(\xi - \eta)$ :
\bea
	p_{\xi} & = & \frac{\partial S(\xi,\eta)}{\partial\xi} = \frac{\partial S(\rho,z)}{\partial\rho}\frac{\partial\rho}{\partial\xi} + \frac{\partial S(\rho,z)}{\partial z}\frac{\partial z}{\partial\eta} = p_{\rho}\frac{\eta}{2\sqrt{\xi\eta}} + \frac{p_{z}}{2} = \frac{r - z}{2\rho}p_{\rho} + \frac{p_{z}}{2} \nonumber \\
	p_{\eta} & = & \frac{\partial S(\xi,\eta)}{\partial\eta} = \frac{\partial S(\rho,z)}{\partial\rho}\frac{\partial\rho}{\partial\eta} + \frac{\partial S(\rho,z)}{\partial z}\frac{\partial z}{\partial\eta} = p_{\rho}\frac{\xi}{2\sqrt{\xi\eta}} - \frac{p_{z}}{2} = \frac{r + z}{2\rho}p_{\rho} + \frac{p_{z}}{2} \nonumber
\eea
Nous avons besoin de l'expression de leur carr\'{e}, soit :
\bea
	p_{\xi}^{2} & = & \frac{(r - z)^{2}}{4\rho^{2}}p_{\rho}^{2} + \frac{p_{z}^{2}}{4} + \frac{(r - z)}{2\rho}p_{\rho}p_{z} = \frac{r^{2}}{4\rho^{2}}p_{\rho}^{2} + \frac{z^{2}}{4\rho^{2}}p_{\rho}^{2} - \frac{rz}{2\rho^{2}}p_{\rho}p_{z} + \frac{p_{z}^{2}}{4} - \frac{r}{2\rho}p_{\rho}p_{z} - \frac{z}{2\rho}p_{\rho}p_{z} \nonumber \\
	p_{\eta}^{2} & = & \frac{(r + z)^{2}}{4\rho^{2}}p_{\rho}^{2} + \frac{p_{z}^{2}}{4} - \frac{(r + z)}{2\rho}p_{\rho}p_{z} = \frac{r^{2}}{4\rho^{2}}p_{\rho}^{2} + \frac{z^{2}}{4\rho^{2}}p_{\rho}^{2} + \frac{rz}{2\rho^{2}}p_{\rho}p_{z} + \frac{p_{z}^{2}}{4} - \frac{r}{2\rho}p_{\rho}p_{z} - \frac{z}{2\rho}p_{\rho}p_{z} \nonumber
\eea
qui nous donne pour :
\benn
	p_{\xi}^{2} - p_{\eta}^{2} = -\frac{rz}{\rho^{2}}p_{\rho}^{2} + \frac{r}{\rho}p_{\rho}p_{z}
\eenn
Ces expressions sont utilis\'{e}es dans le d\'{e}veloppement de la quantité $\beta$ :
\bea
	2\beta & = & \frac{4\rho^{2}}{2r}\left(\frac{r}{\rho}p_{\rho}p_{z} - \frac{rz}{\rho^{2}}p_{\rho}^{2}\right) - \frac{2z}{\rho^{2}}p_{\varphi}^{2} - 2m\alpha\frac{r}{z} - m\rho^{2}F \nonumber \\
	\Leftrightarrow \beta & = & \rho p_{\rho}p_{z} - zp_{\rho}^{2} - \frac{z}{\rho^{2}}p_{\varphi}^{2} - m\alpha\frac{z}{r} - \frac{m}{2}\rho^{2}F \nonumber \\
	\Leftrightarrow \beta & = & -m\left(\frac{\alpha z}{r} + \frac{p_{\rho}}{m}(zp_{\rho} - \rho p_{z}) + \frac{p_{\varphi}^{2}}{m\rho^{2}}z\right) - \frac{m}{2}F\rho^{2} \nonumber
\eea
L'expression entre parenth\`{e}ses correspond \`{a} l'int\'{e}grale premi`{e}re d'un champ purement coulombien d\'{e}fini \`{a} la formule (\ref{EQ:15_17}) et correspondant \`{a} $F = 0$ dans notre exercice.

\subsection{Superposition de deux champs coulombiens}\label{PAR:44_EX2}

Nous consid\'{e}rons le champ potentiel suivant :
\benn
	U = \frac{\alpha_{1}}{r_{1}} + \frac{\alpha_{2}}{r_{2}}
\eenn
de deux centres immobiles distants de $2\sigma$ l'un de l'autre.

\`{A} la vue de l'\'{e}nonc\'{e}, nous utilisons les coordonn\'{e}es elliptiques du paragraphe (\ref{PAR:48_3_3}), soit :
\benn
	r_{1} = \sigma(\xi - \eta)\text{ et }r_{2} = \sigma(\xi + \eta)
\eenn
qui nous permet d\'{e}crire l'\'{e}nergie potentielle ainsi :
\benn
	U = \frac{\alpha_{1}}{\sigma(\xi - \eta)} + \frac{\alpha_{2}}{\sigma(\xi + \eta)} = \frac{(\alpha_{1} + \alpha_{2})\xi + (\alpha_{1} - \alpha_{2})\eta}{\sigma(\xi^{2} - \eta^{2})}
\eenn
En s'appuyant sur la m\'{e}thode de s\'{e}paration des variables pour l'\'{e}nergie potentielle, voir la formule (\ref{EQ:48_2}), nous obtenons les deux fonctions suivantes :
\benn
	a(\xi) = \frac{(\alpha_{1} + \alpha_{2})}{\sigma}\xi\text{ et }b(\eta) = \frac{(\alpha_{1} - \alpha_{2})}{\sigma}\eta
\eenn
L'action du syst\`{e}me s'obtient donc en ins\'{e}rant $a(\xi)$ et $b(\eta)$ dans la formule (\ref{EQ:48_22}). Pour le sens de la quantit\'{e} $\beta$, nous repartons des deux \'{e}quations issues du paragraphe (\ref{PAR:48_3_3}), soit :
\bea
	(\xi^{2} - 1)p_{\xi}^{2} + \frac{1}{\xi^{2} - 1}p_{\varphi}^{2} + 2m\sigma^{2}a(\xi) - 2m\sigma^{2}E\xi^{2} & = & \beta \nonumber \\
	(1 - \eta^{2})p_{\eta}^{2} + \frac{1}{1 - \eta^{2}}p_{\varphi}^{2} + 2m\sigma^{2}b(\eta) + 2m\sigma^{2}E\eta^{2} & = & -\beta \nonumber
\eea
La soustraction de la premi\`{e}re par la seconde donne :
\benn
	2\beta = (\xi^{2} - 1)p_{\xi}^{2} - (1 - \eta^{2})p_{\eta}^{2} + \left(\frac{1}{\xi^{2} - 1} - \frac{1}{1 - \eta^{2}}\right)p_{\varphi}^{2} + 2m\sigma^{2}\left(a(\xi) - b(\eta)\right) - 2m\sigma^{2}E(\xi^{2} + \eta^{2})
\eenn
La conservation de l'\'{e}nergie du syst\`{e}me, $H = E$, permet d'\'{e}crire en prenant appui sur l'expression de la fonction de Hamilont (\ref{EQ:48_20}) :
\bea
	E & = & \frac{1}{2m\sigma^{2}(\xi^{2} - \eta^{2})}\left[(\xi^{2} - 1)p_{\xi}^{2} + (1 - \eta^{2})p_{\eta}^{2} + \left(\frac{1}{\xi^{2} - 1} + \frac{1}{1 - \eta^{2}}\right)p_{\varphi}^{2}\right] \nonumber \\
	& & + \frac{(\alpha_{1} + \alpha_{2})\xi}{\sigma(\xi^{2} - \eta^{2})} + \frac{(\alpha_{1} - \alpha_{2})\eta}{\sigma(\xi^{2} - \eta^{2})} \nonumber
\eea
Cela nous permet de d\'{e}velopper l'expression de la variable $\beta$ :
\bea
	2\beta & = & (\xi^{2} - 1)p_{\xi}^{2} - (1 - \eta^{2})p_{\eta}^{2} + \left(\frac{1}{\xi^{2} - 1} - \frac{1}{1 - \eta^{2})}\right)p_{\varphi}^{2} + 2m\sigma^{2}\left(\frac{(\alpha_{1} + \alpha_{2})\xi}{\sigma} - \frac{(\alpha_{1} - \alpha_{2})\eta}{\sigma}\right) \nonumber \\
	& & -2m\sigma^{2}(\xi^{2} + \eta^{2})\Bigg[\frac{1}{2m\sigma^{2}(\xi^{2} - \eta^{2})}\left[(\xi^{2} - 1)p_{\xi}^{2} + (1 - \eta^{2})p_{\eta}^{2} + \left(\frac{1}{\xi^{2} - 1} + \frac{1}{1 - \eta^{2}}\right)p_{\varphi}^{2}\right] \nonumber \\
	& & + \frac{(\alpha_{1} + \alpha_{2})\xi}{\sigma(\xi^{2} - \eta^{2})} + \frac{(\alpha_{1} - \alpha_{2})\eta}{\sigma(\xi^{2} - \eta^{2})}\Bigg] \nonumber \\
	& = & \left[(\xi^{2} - 1) - \frac{(\xi^{2} + \eta^{2})(\xi^{2} - 1)}{\xi^{2} - \eta^{2})}\right]p_{\xi}^{2} - \left[(1 - \eta^{2}) + \frac{(\xi^{2} + \eta^{2})(1 - \eta^{2})}{\xi^{2} - \eta^{2})}\right]p_{\eta}^{2} \nonumber \\
	& & \left[\frac{1}{\xi^{2} - 1} - \frac{1}{1 - \eta^{2})} - \frac{\xi^{2} + \eta^{2}}{\xi^{2} - \eta^{2}}\left(\frac{1}{\xi^{2} - 1} + \frac{1}{1 - \eta^{2})}\right)\right]p_{\varphi}^{2} \nonumber \\
	& & + 2m\sigma(\alpha_{1} + \alpha_{2})\xi\left(1 - \frac{\xi^{2} + \eta^{2}}{\xi^{2} - \eta^{2}}\right) - 2m\sigma(\alpha_{1} - \alpha_{2})\eta\left(1 + \frac{\xi^{2} + \eta^{2}}{\xi^{2} - \eta^{2}}\right) \nonumber
\eea
Or, nous avons :
\bea
	1 - \frac{\xi^{2} + \eta^{2}}{\xi^{2} - \eta^{2}} & = & -2\frac{\eta^{2}}{\xi^{2} - \eta^{2}} \nonumber \\
	1 + \frac{\xi^{2} + \eta^{2}}{\xi^{2} - \eta^{2}} & = & 2\frac{\xi^{2}}{\xi^{2} - \eta^{2}} \nonumber \\
	\frac{1}{\xi^{2} - 1} - \frac{1}{1 - \eta^{2})} - \frac{\xi^{2} + \eta^{2}}{\xi^{2} - \eta^{2}}\left(\frac{1}{\xi^{2} - 1} + \frac{1}{1 - \eta^{2})}\right) & = & \frac{1 - \eta^{2} - \xi^{2} + 1}{(\xi^{2} - 1)(1 - \eta^{2})} - \frac{\xi^{2} + \eta^{2}}{\xi^{2} - \eta^{2}}\left(\frac{1 - \eta^{2} + \xi^{2} - 1}{(\xi^{2} - 1)(1 - \eta^{2})}\right) \nonumber \\
	& = & \frac{1}{(\xi^{2} - 1)(1 - \eta^{2})}\left[2 - (\xi^{2} + \eta^{2}) - \frac{\xi^{2} + \eta^{2}}{\xi^{2} - \eta^{2}}(\xi^{2} - \eta^{2})\right] \nonumber \\
	& = & 2\frac{\left(1 - (\xi^{2} + \eta^{2})\right)}{(\xi^{2} - 1)(1 - \eta^{2})} \nonumber
\eea
Donc :
\bea
	2\beta & = & -2\frac{(\xi^{2} - 1)}{\xi^{2} - \eta^{2}}\eta^{2}p_{\xi}^{2} - 2\frac{(1 - \eta^{2})}{\xi^{2} - \eta^{2}}\xi^{2}p_{\eta}^{2} - 2\frac{\left(1 - (\xi^{2} + \eta^{2})\right)}{(\xi^{2} - 1)(1 - \eta^{2})}p_{\varphi}^{2} - 4\frac{m\sigma\xi\eta}{\xi^{2} - \eta^{2}}\left(\alpha_{1}(\xi + \eta) + \alpha_{2}(\eta - \xi)\right) \nonumber \\
	& = & -\frac{\xi^{2} - 1}{\xi^{2} - \eta^{2}}\eta^{2}p_{\xi}^{2} - \frac{1 - \eta^{2}}{\xi^{2} - \eta^{2}}\xi^{2}p_{\eta}^{2} - \frac{1 - (\xi^{2} + \eta^{2})}{(\xi^{2} - 1)(1 - \eta^{2})}p_{\varphi}^{2} - 2\frac{m\sigma\xi\eta}{\xi^{2} - \eta^{2}}\left(\alpha_{1}(\xi + \eta) + \alpha_{2}(\eta - \xi)\right) \nonumber
\eea
Nous allons d\'{e}sormais changer de variables pour passer de $(\xi,\eta)$ \`{a} $(\rho,z)$ :
\bea
	p_{\xi} & = & \frac{\partial S(\xi,\eta)}{\partial\xi} = \frac{\partial S(\rho,z)}{\partial\rho}\frac{\partial\rho}{\partial\xi} + \frac{\partial S(\rho,z)}{\partial z}\frac{\partial z}{\partial\xi} = \frac{\sigma(1 - \eta^{2})\xi}{\sqrt{(\xi^{2} - 1)(1 - \eta^{2})}}p_{\rho} + \sigma\eta p_{z} \nonumber \\
	p_{\eta} & = & \frac{\partial S(\xi,\eta)}{\partial\eta} = \frac{\partial S(\rho,z)}{\partial\rho}\frac{\partial\rho}{\partial\eta} + \frac{\partial S(\rho,z)}{\partial z}\frac{\partial z}{\partial\eta} = -\frac{\sigma(\xi^{2} - 1)\eta}{\sqrt{(\xi^{2} - 1)(1 - \eta^{2})}}p_{\rho} + \sigma\xi p_{z} \nonumber
\eea
car $\rho = \sqrt{(\xi^{2} - 1)(1 - \eta^{2})}$ et $z = \sigma\xi\eta$. Le carr\'{e} de ces impulsions sont :
\bea
	p_{\xi}^{2} & = & \frac{\sigma^{2}(1 - \eta^{2})^{2}}{(\xi^{2} - 1)(1 - \eta^{2})}\xi^{2}p_{\rho}^{2} + \sigma^{2}\eta^{2}p_{z}^{2} + 2\frac{\sigma^{2}(1 - \eta^{2})}{\sqrt{(\xi^{2} - 1)(1 - \eta^{2})}}\xi\eta p_{\rho}p_{z} \nonumber \\
	& = & \frac{\sigma^{2}(1 - \eta^{2})}{(\xi^{2} - 1)}\xi^{2}p_{\rho}^{2} + \sigma^{2}\eta^{2}p_{z}^{2} + 2\frac{\sigma^{2}(1 - \eta^{2})}{\sqrt{(\xi^{2} - 1)(1 - \eta^{2})}}\xi\eta p_{\rho}p_{z} \nonumber \\
	p_{\eta}^{2} & = & \frac{\sigma^{2}(\xi^{2} - 1)^{2}}{(\xi^{2} - 1)(1 - \eta^{2})}\eta^{2}p_{\rho}^{2} + \sigma^{2}\xi^{2}p_{z}^{2} - 2\frac{\sigma^{2}(\xi^{2} - 1)}{\sqrt{(\xi^{2} - 1)(1 - \eta^{2})}}\xi\eta p_{\rho}p_{z} \nonumber \\
	& = & \frac{\sigma^{2}(\xi^{2} - 1)}{(1 - \eta^{2})}\eta^{2}p_{\rho}^{2} + \sigma^{2}\xi^{2}p_{z}^{2} - 2\frac{\sigma^{2}(\xi^{2} - 1)}{\sqrt{(\xi^{2} - 1)(1 - \eta^{2})}}\xi\eta p_{\rho}p_{z} \nonumber
\eea
Nous examinons d'abord le terme facteur de $p_{\rho}^{2}$ dans $\beta$ :
\bea
	-\frac{\xi^{2} - 1}{\xi^{2} - \eta^{2}}\eta^{2}\times\frac{\sigma^{2}(1 - \eta^{2})}{\xi^{2} - 1}\xi^{2} - \frac{1 - \eta^{2}}{\xi^{2} - \eta^{2}}\xi^{2}\times\frac{\sigma^{2}(\xi^{2} - 1)}{1 - \eta^{2}}\eta^{2} & = & -\frac{\sigma^{2}(1 - \eta^{2})}{\xi^{2} - \eta^{2}}\xi^{2}\eta^{2} - \frac{\sigma^{2}(\xi^{2} - 1)}{\xi^{2} - \eta^{2}}\xi^{2}\eta^{2} \nonumber \\
	& = & -\frac{\sigma^{2}\xi^{2}\eta^{2}}{\xi^{2} - \eta^{2}}(1 - \eta^{2} + \xi^{2} - 1) = -\sigma^{2}\xi^{2}\eta^{2} \nonumber \\
	& = & -z^{2} \nonumber
\eea
alors que le terme facteur de $p_{z}^{2}$ devient :
\bea
	-\frac{\xi^{2} - 1}{\xi^{2} - \eta^{2}}\eta^{2}\sigma^{2}\eta^{2} - \frac{1 - \eta^{2}}{\xi^{2} - \eta^{2}}\xi^{2}\sigma^{2}\xi^{2} & = & -\frac{\sigma^{2}}{\xi^{2} - \eta^{2}}\left((\xi^{2} - 1)\eta^{4} + (1 - \eta^{2})\xi^{4}\right) \nonumber \\
	& = & -\frac{\sigma^{2}}{\xi^{2} - \eta^{2}}(\xi^{2}\eta^{4} - \eta^{4} + \xi^{4} - \xi^{4}\eta^{2}) = -\frac{\sigma^{2}}{\xi^{2} - \eta^{2}}\left(\xi^{4} - \eta^{4} - \xi^{2}\eta^{2}(\xi^{2} - \eta^{2})\right) \nonumber \\
	& = & -\frac{\sigma^{2}}{\xi^{2} - \eta^{2}}\left((\xi^{2} + \eta^{2})(\xi^{2} - \eta^{2}) - \xi^{2}\eta^{2}(\xi^{2} - \eta^{2})\right) = -\sigma^{2}(\xi^{2} + \eta^{2} - \xi^{2}\eta^{2}) \nonumber \\
	& = & -\rho^{2} - \sigma^{2} - z^{2} + z^{2} = -(\rho^{2} + \sigma^{2}) \nonumber
\eea
car :
\benn
	\rho^{2} = \sigma^{2}(\xi^{2} - 1)(1 - \eta^{2}) = \sigma^{2}(\xi^{2} - \xi^{2}\eta^{2} - 1 + \eta^{2}) = \sigma^{2}(\xi^{2} + \eta^{2}) - \sigma^{2} - z^{2}
\eenn
et celui facteur de $p_{\rho}p_{z}$ :
\bea
	& & -\frac{\xi^{2} - 1}{\xi^{2} - \eta^{2}}\eta^{2}\times\frac{2\sigma^{2}(1 - \eta^{2})}{\sqrt{(\xi^{2} - 1)(1 - \eta^{2})}}\xi\eta - \frac{1 - \eta^{2}}{\xi^{2} - \eta^{2}}\xi^{2}\times\frac{-2\sigma^{2}(\xi^{2} - 1)}{\sqrt{(\xi^{2} - 1)(1 - \eta^{2})}}\xi\eta \nonumber \\
	& = & 2\frac{(\xi^{2} - 1)(1 - \eta^{2})\sigma^{2}}{(\xi^{2} - \eta^{2})\sqrt{(\xi^{2} - 1)(1 - \eta^{2})}}\xi\eta(-\eta^{2} + \xi^{2}) = 2\sigma^{2}\sqrt{(\xi^{2} - 1)(1 - \eta^{2})}\xi\eta \nonumber \\
	& = & 2\rho z \nonumber
\eea
Nous calculons maintenant le terme facteur de $\alpha_{1}$ :
\benn
	-2m\sigma\frac{\xi\eta}{\xi^{2} - \eta^{2}}(\eta + \xi) = -2m\sigma\frac{(r_{2}^{2} - r_{1}^{2})}{4\sigma}\times\frac{\sigma^{2}}{r_{1}r_{2}}\times\frac{r^{2}}{\sigma} = -m\sigma\frac{r_{2}^{2} - r_{1}^{2}}{2r_{1}}
\eenn
car nous avons toujours les formules de transformation suivantes :
\benn
	\begin{cases}
		r_{1} = \sigma(\xi - \eta) \\
		r_{2} = \sigma(\xi + \eta) \\
		r_{1}r_{2} = \sigma^{2}(\xi^{2} - \eta^{2}) \\
		\xi = \frac{r_{1} + r_{2}}{2\sigma} \\
		\eta = \frac{r_{2} - r_{1}}{2\sigma} \\
		\rho^{2} = \sigma^{2}(\xi^{2} - 1)(1 - \eta^{2}) = \sigma^{2}(\xi^{2} - \eta^{2}) - \sigma^{2} - z^{2} \\
		z = \sigma\xi\eta = \frac{r_{2}^{2} - r_{1}}{4\sigma}
	\end{cases}
\eenn
Et le terme facteur de $\alpha_{2}$ est :
\benn
	-2m\sigma\frac{\xi\eta}{\xi^{2} - \eta^{2}}(\eta - \xi) = 2m\sigma\frac{\xi\eta}{\xi^{2} - \eta^{2}}(\xi - \eta) = 2m\sigma\frac{(r_{2}^{2} - r_{1}^{2})}{4\sigma}\times\frac{\sigma^{2}}{r_{1}r_{2}}\times\frac{r_{1}}{\sigma} = m\sigma\frac{r_{2}^{2} - r_{1}^{2}}{2r_{2}}
\eenn
Enfin le terme facteur de $p_{\varphi}^{2}$ est :
\benn
	-\frac{\left[1 - (\xi^{2} + \eta^{2})\right]}{(\xi^{2} - 1)(1 - \eta^{2})}= -\frac{\left[1 - \left(\frac{\rho^{2}}{\sigma^{2}} + 1 + \frac{z^{2}}{\sigma^{2}}\right)\right]}{\rho^{2}/\sigma^{2}} = -\frac{\sigma^{2} - \rho^{2} - \sigma^{2} - z^{2}}{\rho^{2}} = 1 + \frac{z^{2}}{\rho^{2}}
\eenn
Finalement, nous pouvons formuler la variable $\beta$ ainsi :
\benn
	\beta = -z^{2}p_{\rho}^{2} - (\rho^{2} + z^{2})p_{z}^{2} + 2\rho zp_{\rho}p_{z} + \left(1 + \frac{z^{2}}{\rho^{2}}\right)p_{\varphi}^{2} - m\sigma\left[\frac{r_{2}^{2} - r_{1}^{2}}{2r_{1}}\alpha_{1} - \frac{r_{2}^{2} - r_{1}^{2}}{2r_{2}}\alpha_{2}\right]
\eenn
Cette expression est l\'{e}g\`{e}rement diff\'{e}rente \`{a} celle exprim\'{e}e dans le livre.

\section{Variables canoniques}\label{PAR:50_EX}

Nous devons \'{e}crire l'\'{e}quation du mouvement d'un oscillateur harmonique \`{a} fr\'{e}quence d\'{e}pendant du temps en utilisant les variables canoniques comme pr\'{e}sent\'{e}es dans le paragraphe (\ref{PAR:50}) et en particulier la d\'{e}finition des \'{e}quations de Hamilton dans les formules (\ref{EQ:50_4}).

La fonction de Hamilton du syst\`{e}me est celle obtenue dans le r\'{e}sultat (\ref{EQ:49_11}) :
\benn
	H = \frac{p^{2}}{2m} + \frac{m\omega^{2}}{2}q^{2}
\eenn
En application de la conservation de l'\'{e}nergie le long de la trajectoire :
\benn
	H = E \Leftrightarrow \frac{p^{2}}{2mE} + \frac{m\omega^{2}}{2E}q^{2} = 1
\eenn
qui est l'\'{e}quation d'une ellipse dans l'espace des phases. Son \'{e}quation peut aussi s'\'{e}crire avec la variable d'angle $w$ en s'appuyant sur $\cos^{2}w + \sin^{2}w = 1$. Nous obtenons les \'{e}galit\'{e}s suivantes :
\bea
	p^{2} & = & 2mE\cos^{2}(w) \Leftrightarrow p = \pm\sqrt{2mE}\cos(w) = \pm\sqrt{2m\omega I}\cos(w) \nonumber \\
	q^{2} & = & \frac{2E}{m\omega^{2}}\sin^{2}(w) \Leftrightarrow q = \pm\sqrt{\frac{2E}{m\omega^{2}}}\sin(w) = \pm\sqrt{\frac{2I}{m\omega}}\sin(w)
\eea
car $E = I\omega$ dans la relation (\ref{EQ:49_12}).

Nous notons que la fr\'{e}quence $\omega(\mathrm{t})$ joue le r\^{o}le de la variable $\lambda$ de la formule (\ref{EQ:50_2}). Par d\'{e}finition (\ref{EQ:50_1}), $S_{0} = \int{p\mathrm{d}q}$ or :
\bea
	\mathrm{d}q & = & \left(\frac{\partial q}{\partial\omega}\right)_{I,w}\mathrm{d}\omega + \left(\frac{\partial q}{\partial I}\right)_{q,I}\mathrm{d}I + \left(\frac{\partial q}{\partial w}\right)_{I,\omega}\mathrm{d}w \nonumber \\
	& = & \sqrt{\frac{2I}{m}}\cdot\left(-\frac{1}{2}\right)\omega^{-3/2}\sin(w)\mathrm{d}\omega + \sqrt{\frac{2}{m\omega}}\cdot\frac{1}{2}I^{-1/2}\sin(w)\mathrm{d}I + \sqrt{\frac{2I}{m\omega}}\cos(w)\mathrm{d}w \nonumber \\
	& = & -\sqrt{\frac{I}{2m}}\omega^{-3/2}\sin(w)\mathrm{d}\omega + \sqrt{\frac{1}{2m\omega}}I^{-1/2}\sin(w)\mathrm{d}I + \sqrt{\frac{2I}{m\omega}}\cos(w)\mathrm{d}w \nonumber
\eea
donc la quantit\'{e} $p\mathrm{d}q$ vaut :
\bea
	p\mathrm{d}q & = & -\sqrt{2m\omega I}\sqrt{\frac{I}{2m}}w^{-3/2}\cos(w)\sin(w)\mathrm{d}\omega + \sqrt{2m\omega I}\sqrt{\frac{1}{2m\omega}}I^{-1/2}\cos(w)\sin(w)\mathrm{d}I \nonumber \\
	& & + \sqrt{2m\omega I}\sqrt{\frac{2I}{m\omega}}\cos^{2}(w)\mathrm{d}w \nonumber \\
	& = & -\frac{I}{\omega}\cos(w)\sin(w)\mathrm{d}\omega + \cos(w)\sin(w)\mathrm{d}I + 2I\cos^{2}(w)\mathrm{d}w \nonumber \\
	& = & \left(\mathrm{d}I - \frac{I}{\omega}\mathrm{d}\omega\right)\cos(w)\sin(w) + 2I\cos^{2}(w)\mathrm{d}w \nonumber
\eea
or la relation (\ref{EQ:49_12}) donne :
\benn
	I = \frac{E}{\omega} \Leftrightarrow \mathrm{d}I = -\frac{E}{\omega^{2}}\mathrm{d}\omega = -\frac{I}{\omega}\mathrm{d}\omega
\eenn
donc :
\benn
	p\mathrm{d}q = 2I\cos^{2}(w)\mathrm{d}w = \left(\frac{\partial q}{\partial w}\right)_{\omega,I}\mathrm{d}q
\eenn
et finalement :
\benn
	S_{0} = 2I\int{cos^{2}(w)\mathrm{d}w}
\eenn
Par d\'{e}finition (\ref{EQ:50_9}), $\Lambda= \left(\frac{\partial S_{0}}{\partial\omega}\right)_{q,I}$ que nous pouvons d\'{e}composer pour son calcul en :
\benn
	\Lambda = \left(\frac{\partial S_{0}}{\partial w}\right)_{I}\cdot\left(\frac{\partial w}{\partial\omega}\right)_{q}
\eenn
or $\frac{\partial S_{0}}{\partial w} = 2I\cos^{2}(w)$ et comme $q = \sqrt{\frac{2I}{m\omega}}\sin(w)$, alors :
\bea
	\frac{\partial w}{\partial\omega} & = & \frac{\partial\sin(w)}{\partial\omega}\cdot\frac{\partial w}{\partial\sin(w)} = \frac{\partial\sin(w)}{\partial\omega}\cdot\left(\frac{\partial\sin(w)}{\partial w}\right)^{-1} = \frac{1}{2}\sqrt{\frac{m}{2I\omega}}q\times\frac{1}{\cos(w)} \nonumber \\
	& = & \frac{1}{2}\sqrt{\frac{m}{2I\omega}}\sqrt{\frac{2I}{m\omega}}\frac{\sin(w)}{\cos(w)} = \frac{1}{2\omega}\frac{\sin(w)}{\cos(w)} \nonumber
\eea
Alors, la quantit\'{e} $\Lambda$ vaut :
\benn
	\Lambda = 2I\cos^{2}(w)\times\frac{1}{2\omega}\frac{\sin(w)}{\cos(w)} = \frac{I}{\omega}\cos(w)\sin(w) = \frac{I}{2\omega}\sin(2w)
\eenn
par application de la formule de Simpson : $\sin(a + b) = \sin(a)\cos(b) + \cos(a)\sin(b)$. Les \'{e}quations de Hamilton sont donn\'{e}es pour les variables canoniques par les formules (\ref{EQ:50_10}) et (\ref{EQ:50_11}), soient :
\bea
	\dot{I} & = & -\left(\frac{\partial\Lambda}{\partial w}\right)_{I,\omega}\dot{\omega} = -I\frac{\dot{\omega}}{\omega}\cos(2w) \nonumber \\
	\dot{\omega} & = & \omega + \left(\frac{\partial\Lambda}{\partial I}\right)_{w,\omega}\dot{\omega} = \omega + \frac{\dot{\omega}}{2\omega}\sin(2w) \nonumber
\eea

\section{Exactitude de la conservation d'invariant adiabatique}\label{PAR:51_EX}

Nous devons ici \'{e}valuer $\Delta I$ dans le cas d'oscillations harmoniques dont la fr\'{e}quence varie lentement suivant la loi :
\benn
	\omega^{2} = \omega_{0}^{2}\frac{1 + ae^{\alpha\mathrm{t}}}{1 + e^{\alpha\mathrm{t}}}
\eenn
avec $a > 0$ et $\alpha \ll \omega_{0}$/ Pour $\mathrm{t} = 0$, $\omega_{-} = \omega_{0}$ et pour $\mathrm{t} \rightarrow +\infty$, $\omega^{2} \rightarrow \omega_{0}^{2}\frac{ae^{\alpha\mathrm{t}}}{e^{\alpha\mathrm{t}}} = a\omega_{0}^{2}$ donc $\omega_{+} = \sqrt{a}\omega_{0}$.
Pour obtenir $\Delta I$ \`{a} partir de l'approximation (\ref{EQ:51_10}), il nous faut d'abord calculer le rapport $\dot{\lambda}/\omega$ sachant que dans le cas pr\'{e}sent $\lambda$ est la fr\'{e}quence $\omega$ elle-m\^{e}me. Or :
\bea
	\dot{\lambda} = \dot{\omega} & = & \omega_{0}\sqrt{\frac{1 + e^{\alpha\mathrm{t}}}{1 + ae^{\alpha\mathrm{t}}}}\frac{\left(a\alpha e^{\alpha\mathrm{t}}\left(1 + e^{\alpha\mathrm{t}}\right) - \left(1 + ae^{\alpha\mathrm{t}}\right)\alpha e^{\alpha\mathrm{t}}\right)}{\left(1 + e^{\alpha\mathrm{t}}\right)^{2}} \nonumber \\
	& = & \frac{\omega}{2}\frac{1 + e^{\alpha\mathrm{t}}}{1 + ae^{\alpha\mathrm{t}}}\frac{a\alpha e^{\alpha\mathrm{t}} + a\alpha e^{2\alpha\mathrm{t}} - \alpha e^{\alpha\mathrm{t}} - a\alpha e^{2\alpha\mathrm{t}}}{\left(1 + e^{\alpha\mathrm{t}}\right)^{2}} \nonumber \\
	& = & \frac{\omega\alpha}{2}\frac{(a - 1)e^{\alpha\mathrm{t}}}{\left(1 + ae^{\alpha\mathrm{t}}\right)\left(1 + e^{\alpha\mathrm{t}}\right)} \nonumber \\
	\Leftrightarrow \frac{\dot{\lambda}}{\omega} & = & \frac{\alpha}{2}\frac{(a - 1)}{\left(a + e^{-\alpha\mathrm{t}}\right)\left(1 + e^{-\alpha\mathrm{t}}\right)} \nonumber
\eea
Cette fonction obtenue a deux p\^{o}les tels que $e^{-\alpha\mathrm{t}} = -1$ et $e^{-\alpha\mathrm{t}} = -a$.
(Voir \emph{https://github.com/SMajerowicz/Physics/issues/31} pour le calcul de $\int w\mathrm{dt}$)
Le calcul de l'int\'{e}grale $\int\omega\mathrm{dt}$ montre que la valeur minimale pour $\Im(w_{0})$ est obtenue pour $\alpha\mathrm{t}_{0} = -\ln(-a)$ et elle vaut :
\benn
	\Im(w_{0}) =
	\begin{cases}
		\omega_{0}\pi/\alpha\text{ avec }a > 1 \\
		\sqrt{a}\omega_{0}\pi/\alpha\text{ avec }a < 1
	\end{cases}
\eenn
L'exercice (\ref{PAR:50_EX}) a permis de montrer que pour un oscillateur harmonique $\Lambda \propto \sin(2w)$. Ainsi la s\'{e}rie de Fourier de $\Lambda$ en (\ref{EQ:51_3}) s'arr\^{e}te exactement \`{a} $n = \pm 2$ et donc $\Delta I \propto e^{-2\Im(w_{0})}$ en adaptant (\ref{EQ:51_6}).

\section{Mouvement quasi-p\'{e}riodique}\label{PAR:52_EX}
